% Options for packages loaded elsewhere
% Options for packages loaded elsewhere
\PassOptionsToPackage{unicode}{hyperref}
\PassOptionsToPackage{hyphens}{url}
\PassOptionsToPackage{dvipsnames,svgnames,x11names}{xcolor}
%
\documentclass[
  spanish,
  letterpaper,
  DIV=11,
  numbers=noendperiod]{scrartcl}
\usepackage{xcolor}
\usepackage{amsmath,amssymb}
\setcounter{secnumdepth}{-\maxdimen} % remove section numbering
\usepackage{iftex}
\ifPDFTeX
  \usepackage[T1]{fontenc}
  \usepackage[utf8]{inputenc}
  \usepackage{textcomp} % provide euro and other symbols
\else % if luatex or xetex
  \usepackage{unicode-math} % this also loads fontspec
  \defaultfontfeatures{Scale=MatchLowercase}
  \defaultfontfeatures[\rmfamily]{Ligatures=TeX,Scale=1}
\fi
\usepackage{lmodern}
\ifPDFTeX\else
  % xetex/luatex font selection
\fi
% Use upquote if available, for straight quotes in verbatim environments
\IfFileExists{upquote.sty}{\usepackage{upquote}}{}
\IfFileExists{microtype.sty}{% use microtype if available
  \usepackage[]{microtype}
  \UseMicrotypeSet[protrusion]{basicmath} % disable protrusion for tt fonts
}{}
\makeatletter
\@ifundefined{KOMAClassName}{% if non-KOMA class
  \IfFileExists{parskip.sty}{%
    \usepackage{parskip}
  }{% else
    \setlength{\parindent}{0pt}
    \setlength{\parskip}{6pt plus 2pt minus 1pt}}
}{% if KOMA class
  \KOMAoptions{parskip=half}}
\makeatother
% Make \paragraph and \subparagraph free-standing
\makeatletter
\ifx\paragraph\undefined\else
  \let\oldparagraph\paragraph
  \renewcommand{\paragraph}{
    \@ifstar
      \xxxParagraphStar
      \xxxParagraphNoStar
  }
  \newcommand{\xxxParagraphStar}[1]{\oldparagraph*{#1}\mbox{}}
  \newcommand{\xxxParagraphNoStar}[1]{\oldparagraph{#1}\mbox{}}
\fi
\ifx\subparagraph\undefined\else
  \let\oldsubparagraph\subparagraph
  \renewcommand{\subparagraph}{
    \@ifstar
      \xxxSubParagraphStar
      \xxxSubParagraphNoStar
  }
  \newcommand{\xxxSubParagraphStar}[1]{\oldsubparagraph*{#1}\mbox{}}
  \newcommand{\xxxSubParagraphNoStar}[1]{\oldsubparagraph{#1}\mbox{}}
\fi
\makeatother


\usepackage{longtable,booktabs,array}
\newcounter{none} % for unnumbered tables
\usepackage{calc} % for calculating minipage widths
% Correct order of tables after \paragraph or \subparagraph
\usepackage{etoolbox}
\makeatletter
\patchcmd\longtable{\par}{\if@noskipsec\mbox{}\fi\par}{}{}
\makeatother
% Allow footnotes in longtable head/foot
\IfFileExists{footnotehyper.sty}{\usepackage{footnotehyper}}{\usepackage{footnote}}
\makesavenoteenv{longtable}
\usepackage{graphicx}
\makeatletter
\newsavebox\pandoc@box
\newcommand*\pandocbounded[1]{% scales image to fit in text height/width
  \sbox\pandoc@box{#1}%
  \Gscale@div\@tempa{\textheight}{\dimexpr\ht\pandoc@box+\dp\pandoc@box\relax}%
  \Gscale@div\@tempb{\linewidth}{\wd\pandoc@box}%
  \ifdim\@tempb\p@<\@tempa\p@\let\@tempa\@tempb\fi% select the smaller of both
  \ifdim\@tempa\p@<\p@\scalebox{\@tempa}{\usebox\pandoc@box}%
  \else\usebox{\pandoc@box}%
  \fi%
}
% Set default figure placement to htbp
\def\fps@figure{htbp}
\makeatother



\ifLuaTeX
\usepackage[bidi=basic,shorthands=off]{babel}
\else
\usepackage[bidi=default,shorthands=off]{babel}
\fi
\ifLuaTeX
  \usepackage{selnolig} % disable illegal ligatures
\fi


\setlength{\emergencystretch}{3em} % prevent overfull lines

\providecommand{\tightlist}{%
  \setlength{\itemsep}{0pt}\setlength{\parskip}{0pt}}



 


\KOMAoption{captions}{tableheading}
\makeatletter
\@ifpackageloaded{caption}{}{\usepackage{caption}}
\AtBeginDocument{%
\ifdefined\contentsname
  \renewcommand*\contentsname{Tabla de contenidos}
\else
  \newcommand\contentsname{Tabla de contenidos}
\fi
\ifdefined\listfigurename
  \renewcommand*\listfigurename{Listado de Figuras}
\else
  \newcommand\listfigurename{Listado de Figuras}
\fi
\ifdefined\listtablename
  \renewcommand*\listtablename{Listado de Tablas}
\else
  \newcommand\listtablename{Listado de Tablas}
\fi
\ifdefined\figurename
  \renewcommand*\figurename{Figura}
\else
  \newcommand\figurename{Figura}
\fi
\ifdefined\tablename
  \renewcommand*\tablename{Tabla}
\else
  \newcommand\tablename{Tabla}
\fi
}
\@ifpackageloaded{float}{}{\usepackage{float}}
\floatstyle{ruled}
\@ifundefined{c@chapter}{\newfloat{codelisting}{h}{lop}}{\newfloat{codelisting}{h}{lop}[chapter]}
\floatname{codelisting}{Listado}
\newcommand*\listoflistings{\listof{codelisting}{Listado de Listados}}
\makeatother
\makeatletter
\makeatother
\makeatletter
\@ifpackageloaded{caption}{}{\usepackage{caption}}
\@ifpackageloaded{subcaption}{}{\usepackage{subcaption}}
\makeatother
\usepackage{bookmark}
\IfFileExists{xurl.sty}{\usepackage{xurl}}{} % add URL line breaks if available
\urlstyle{same}
% fallback for those not using the hyperref driver hyperxmp:
\makeatletter
\@ifundefined{xmpquote}{\newcommand{\xmpquote}[1]{#1}}{}
\makeatother
\hypersetup{
  pdftitle={Contenido de nitrógeno en plantas de trébol},
  pdflang={es},
  colorlinks=true,
  linkcolor={blue},
  filecolor={Maroon},
  citecolor={Blue},
  urlcolor={Blue},
  pdfcreator={LaTeX via pandoc}}


\title{Contenido de nitrógeno en plantas de trébol}
\usepackage{etoolbox}
\makeatletter
\providecommand{\subtitle}[1]{% add subtitle to \maketitle
  \apptocmd{\@title}{\par {\large #1 \par}}{}{}
}
\makeatother
\subtitle{Análisis de un diseño completamente al azar}
\author{}
\date{}
\begin{document}
\maketitle


\section{Materiales y métodos}\label{materiales-y-muxe9todos}

\subsection{Diseño experimental}\label{diseuxf1o-experimental}

Se utilizó un diseño completamente al azar con 6 tratamientos y 5
repeticiones por tratamiento, para un total de 30 unidades
experimentales. Se evaluaron tratamientos de inoculación con cepas de
\emph{Rhizobium}, identificados como 3D0K1, 3D0K13, 3D0K4, 3D0K5, 3D0K7
y Compuesto. La unidad experimental fue una matera con una planta de
trébol, inoculada a los 30 días después de la siembra. La variable
respuesta fue el contenido de nitrógeno en el follaje, expresado en
miligramos.

La descripción del procedimiento de cuantificación y su referencia
bibliográfica deben completarse con el protocolo efectivamente
utilizado; la referencia incompleta del documento base no permite
identificar el método.

\subsection{Análisis estadístico}\label{anuxe1lisis-estaduxedstico}

El análisis se realizó en el ambiente R. Se ajustó un modelo lineal con
el tratamiento como efecto fijo para evaluar las diferencias en el
contenido de nitrógeno. Se utilizó la prueba de Levene centrada en la
mediana (variante de Brown--Forsythe) para examinar la homogeneidad de
varianzas y la prueba de Shapiro--Wilk para examinar la normalidad de
los residuos. La independencia de los errores se sustenta en la
aleatorización y en la identificación de unidades experimentales
distintas; no se demuestra mediante estas pruebas.

El efecto de los tratamientos se evaluó mediante el análisis de
varianza. Las medias estimadas por el modelo se obtuvieron con el
paquete \textbf{emmeans} y la comparación múltiple se realizó mediante
la prueba de Tukey, con un nivel de significancia del 5 \%. Se utilizó
\textbf{multcomp}, con apoyo de \textbf{multcompView}, para representar
las comparaciones mediante letras, y \textbf{car} para la prueba de
Levene. Se calcularon intervalos de confianza del 95 \% para las medias
estimadas por el modelo. Las letras corresponden a comparaciones de
Tukey y no deben confundirse con el ajuste de los intervalos para cada
media.

\section{Resultados y discusión}\label{resultados-y-discusiuxf3n}

Según la prueba de Levene, no se rechazó la hipótesis nula de
homogeneidad de varianzas (p = 0,4794). Este resultado no demuestra que
las varianzas sean iguales. La prueba de Shapiro--Wilk aplicada a los
residuos señaló evidencia contra la normalidad (p = 0,0149).

El análisis de varianza detectó diferencias estadísticamente
significativas entre las medias de los tratamientos (F = 14,3705; gl = 5
y 24; p = 1,485e-06). Esto indica que al menos una media difiere de las
demás.

Los dos mayores valores estimados de contenido de nitrógeno
correspondieron a 3D0K1 y 3D0K5 (28,82 y 23,98 mg, respectivamente;
Figura~\ref{fig-nitrogeno}). Este orden describe las estimaciones
puntuales y no implica por sí solo que ambos superen a todos los demás
tratamientos.

Al comparar 3D0K1 con los otros tratamientos, la prueba de Tukey detectó
un contenido de nitrógeno mayor que en 3D0K7, Compuesto, 3D0K4 y 3D0K13.
No se detectaron diferencias significativas frente a 3D0K5; ello no
demuestra equivalencia entre tratamientos. Las comparaciones se
interpretan bajo el modelo de varianza común.

Dado que al menos un diagnóstico señaló una desviación de los supuestos,
la ANOVA y las comparaciones de Tukey se presentan como resultados del
modelo clásico ajustado y deben interpretarse con cautela. Es necesario
revisar los gráficos de residuos y considerar un análisis de
sensibilidad adecuado antes de formular conclusiones definitivas. El
contraste de Levene no significativo, si corresponde, no compensa una
desviación de normalidad.

\begin{figure}

\centering{

\includegraphics[width=1\linewidth,height=\textheight,keepaspectratio]{../salidas/graficos/lineas_clasico_homogeneo.jpg}

}

\caption{\label{fig-nitrogeno}Contenido de nitrógeno en el follaje de
plantas de trébol según el tratamiento. Los valores estimados por el
modelo se resaltan en negro, la mediana en naranja y las observaciones
originales en azul, según el gráfico producido por el programa adjunto.
Las barras de error representan los límites de confianza del 95 \%
guardados por emmeans. Tratamientos que no comparten letras presentan
diferencias significativas según Tukey al 5 \%. Compartir una letra
indica que no se detectó una diferencia; no demuestra equivalencia.}

\end{figure}%

\subsection{Medias estimadas y comparación de
tratamientos}\label{medias-estimadas-y-comparaciuxf3n-de-tratamientos}

{

\begin{longtable}[]{@{}
  >{\raggedright\arraybackslash}p{(\linewidth - 10\tabcolsep) * \real{0.1304}}
  >{\raggedleft\arraybackslash}p{(\linewidth - 10\tabcolsep) * \real{0.2174}}
  >{\raggedleft\arraybackslash}p{(\linewidth - 10\tabcolsep) * \real{0.1630}}
  >{\raggedleft\arraybackslash}p{(\linewidth - 10\tabcolsep) * \real{0.1739}}
  >{\raggedleft\arraybackslash}p{(\linewidth - 10\tabcolsep) * \real{0.1739}}
  >{\centering\arraybackslash}p{(\linewidth - 10\tabcolsep) * \real{0.1413}}@{}}

\caption{\label{tbl-medias}Medias estimadas del contenido de nitrógeno,
error estándar y límites de confianza del 95 \%. Las letras corresponden
a la comparación múltiple de Tukey.}

\tabularnewline

\toprule\noalign{}
\begin{minipage}[b]{\linewidth}\raggedright
Tratamiento
\end{minipage} & \begin{minipage}[b]{\linewidth}\raggedleft
Media estimada (mg)
\end{minipage} & \begin{minipage}[b]{\linewidth}\raggedleft
Error estándar
\end{minipage} & \begin{minipage}[b]{\linewidth}\raggedleft
Límite inferior
\end{minipage} & \begin{minipage}[b]{\linewidth}\raggedleft
Límite superior
\end{minipage} & \begin{minipage}[b]{\linewidth}\centering
Grupo Tukey
\end{minipage} \\
\midrule\noalign{}
\endhead
\bottomrule\noalign{}
\endlastfoot
3D0K1 & 28,82 & 1,54 & 24,42 & 33,22 & a \\
3D0K5 & 23,98 & 1,54 & 19,58 & 28,38 & ab \\
3D0K7 & 19,92 & 1,54 & 15,52 & 24,32 & bc \\
Compuesto & 18,70 & 1,54 & 14,30 & 23,10 & bc \\
3D0K4 & 14,64 & 1,54 & 10,24 & 19,04 & c \\
3D0K13 & 13,26 & 1,54 & 8,86 & 17,66 & c \\

\end{longtable}

}

Los límites anteriores corresponden a un nivel de confianza del 95 \% y
conservan el ajuste de la salida original de emmeans; no se interpretan
como intervalos de Tukey para diferencias de medias. La ausencia de
diferencias significativas entre dos tratamientos no permite concluir
que sus efectos sean equivalentes.

\subsection{Resumen de las pruebas}\label{resumen-de-las-pruebas}

{

\begin{longtable}[]{@{}lr@{}}

\caption{\label{tbl-pruebas}Valores p de las pruebas de homogeneidad de
varianzas, normalidad de los residuos y efecto del tratamiento.}

\tabularnewline

\toprule\noalign{}
Prueba & Valor p \\
\midrule\noalign{}
\endhead
\bottomrule\noalign{}
\endlastfoot
Levene & p = 0,4794 \\
Shapiro--Wilk de los residuos & p = 0,0149 \\
ANOVA del tratamiento & p = 1,485e-06 \\

\end{longtable}

}




\end{document}
